% from ega3-1-fr.tex
\oldpage[III]{90}
\begin{proposition}[1.4.5]
\label{III.1.4.5-fr}
Soient $X$ un schéma quasi-compact, $\mathcal L$ un $\mathcal O_X$-Module
inversible, et considérons l'anneau gradué
$A_\bullet=\Gamma_\bullet(\mathcal L)$ (0\textsubscript{I}, 5.4.6)~; alors
\[
H^\bullet(\mathcal F,\mathcal L)
=\bigoplus_{n\in\mathbb Z}H^\bullet(X,\mathcal F\otimes\mathcal L^{\otimes n})
\]
est un $A_\bullet$-module gradué, et pour tout $f\in A_n$, on a un
isomorphisme canonique
\begin{equation}
\label{III.1.4.5.1-fr}
\tag{1.4.5.1}
H^\bullet(X_f,\mathcal F)
\xrightarrow{\sim}(H^\bullet(\mathcal F,\mathcal L))_{(f)}
\end{equation}
de $(A_\bullet)_{(f)}$-modules.
\end{proposition}

\begin{proof}
Comme $X$ est un schéma quasi-compact, on peut calculer la cohomologie de
tous les $\mathcal O_X$-Modules
$\mathcal F\otimes\mathcal L^{\otimes n}$ à l'aide d'un même recouvrement
fini $\mathfrak U=(U_i)$ par des ouverts affines tels que la restriction
$\mathcal L\vert U_i$ soit isomorphe à $\mathcal O_X\vert U_i$ pour chaque
$i$ (\hyperref[III.1.4.1-fr]{1.4.1}). Il est alors immédiat que les
$U_i\cap X_f$ sont des ouverts affines (I, 1.3.6), et l'on peut donc aussi
calculer la cohomologie
$H^\bullet(X_f,\mathcal F\otimes\mathcal L^{\otimes n})$ à l'aide du
recouvrement $\mathfrak U\vert X_f=(U_i\cap X_f)$
(\hyperref[III.1.4.1-fr]{1.4.1}). Il est immédiat que pour tout $f\in A_n$,
la multiplication par $f$ définit un homomorphisme
\[
C^\bullet(\mathfrak U,\mathcal F\otimes\mathcal L^{\otimes m})
\longrightarrow
C^\bullet(\mathfrak U,\mathcal F\otimes\mathcal L^{\otimes(m+n)}),
\]
d'où un homomorphisme
\[
H^\bullet(\mathfrak U,\mathcal F\otimes\mathcal L^{\otimes m})
\longrightarrow
H^\bullet(\mathfrak U,\mathcal F\otimes\mathcal L^{\otimes(m+n)}),
\]
ce qui établit la première assertion. D'autre part, pour $f\in A_n$ donné,
il résulte de (I, 9.3.2) que l'on a un isomorphisme de complexes de
$(A_\bullet)_{(f)}$-modules
\[
C^\bullet(\mathfrak U\vert X_f,\mathcal F)
\xrightarrow{\sim}
\left(C^\bullet\left(\mathfrak U,
\bigoplus_{n\in\mathbb Z}\mathcal F\otimes\mathcal L^{\otimes n}
\right)\right)_{(f)}
\]
compte tenu de (I, 1.3.9, (ii)). Passant à la cohomologie de ces deux
complexes, on en déduit l'isomorphisme
(\hyperref[III.1.4.5.1-fr]{1.4.5.1}), en se souvenant que le foncteur
$M\mapsto M_{(f)}$ est exact dans la catégorie des $A_\bullet$-modules
gradués.
\end{proof}

\begin{corollary}[1.4.6]
\label{III.1.4.6-fr}
On suppose vérifiées les hypothèses de
(\hyperref[III.1.4.5-fr]{1.4.5}) et on suppose en outre que
$\mathcal L=\mathcal O_X$. Si on pose $A=\Gamma(X,\mathcal O_X)$, alors,
pour tout $f\in A$, on a un isomorphisme canonique
\[
H^\bullet(X_f,\mathcal F)
\xrightarrow{\sim}(H^\bullet(X,\mathcal F))_f
\]
de $A_f$-modules.
\end{corollary}

\begin{corollary}[1.4.7]
\label{III.1.4.7-fr}
Soient $X$ un schéma quasi-compact, $f$ un élément de
$\Gamma(X,\mathcal O_X)$.
\begin{enumerate}
\item[(i)] Supposons l'ouvert $X_f$ affine. Alors, pour tout
$\mathcal O_X$-Module quasi-cohérent $\mathcal F$, tout $i>0$ et tout
$\xi\in H^i(X,\mathcal F)$, il existe un entier $n>0$ tel que $f^n\xi=0$.
\item[(ii)] Inversement, supposons que $X_f$ soit quasi-compact et que pour
tout faisceau quasi-cohérent $\mathcal J$ d'idéaux de $\mathcal O_X$ et tout
$\zeta\in H^1(X,\mathcal J)$, il existe $n>0$ tel que $f^n\zeta=0$. Alors
$X_f$ est affine.
\end{enumerate}
\end{corollary}

\begin{proof}
\begin{enumerate}
\item[(i)] Si $X_f$ est affine, on a $H^i(X_f,\mathcal F)=0$ pour tout
$i>0$ (\hyperref[III.1.3.1-fr]{1.3.1}), donc l'assertion résulte directement
de (\hyperref[III.1.4.6-fr]{1.4.6}).
\item[(ii)] En vertu du critère de Serre (II, 5.2.1), il suffit de prouver
que pour tout Idéal quasi-cohérent $\mathcal K$ de
$\mathcal O_X\vert X_f$, on a $H^1(X_f,\mathcal K)=0$. Comme $X_f$ est un
ouvert quasi-compact dans un schéma quasi-compact $X$, il existe un Idéal
quasi-cohérent $\mathcal J$ de $\mathcal O_X$ tel que
$\mathcal K=\mathcal J\vert X_f$ (I, 9.4.2). En vertu de
(\hyperref[III.1.4.6-fr]{1.4.6}), on a
$H^1(X_f,\mathcal K)=(H^1(X,\mathcal J))_f$, et l'hypothèse entraîne que le
second membre est nul, d'où la conclusion.
\end{enumerate}
\end{proof}

\begin{remark}[1.4.8]
\label{III.1.4.8-fr}
On notera que (\hyperref[III.1.4.7-fr]{1.4.7}, (i)) redonne une
démonstration plus simple de la relation (II, 4.5.13.2).
\end{remark}

